% ===== PRÉAMBULE CANONIQUE — à coller VERBATIM en tête de CHAQUE .tex =====
\documentclass[11pt,a4paper]{article}
\usepackage[utf8]{inputenc}\usepackage[T1]{fontenc}\usepackage{lmodern}
\usepackage[french]{babel}
\usepackage{amsmath,amssymb,mathtools}\usepackage{icomma}
\usepackage{siunitx}\sisetup{locale=FR}  % \qty{}{}, \si{}, \ang{}
\usepackage{xcolor}\usepackage{colortbl}
\usepackage{tikz}\usepackage{pgfplots}\pgfplotsset{compat=1.18}
\usepackage[siunitx,europeanresistors]{circuitikz} % circuits (élec.)
\usepackage[most]{tcolorbox}\usepackage{enumitem}\usepackage{array,booktabs}
\usepackage[a4paper,margin=2.2cm]{geometry}
\usetikzlibrary{arrows.meta,calc,angles,quotes,positioning,patterns,%
  backgrounds,shapes.geometric,decorations.pathmorphing,decorations.markings,intersections}
\definecolor{primary}{RGB}{222,49,99}\definecolor{primarydark}{RGB}{150,22,55}
\definecolor{defcol}{RGB}{0,114,178}\definecolor{methcol}{RGB}{52,92,148}
\definecolor{excol}{RGB}{0,135,90}\definecolor{attncol}{RGB}{200,40,55}
\tcbset{boxbase/.style={enhanced,breakable,boxrule=0.9pt,arc=2.5pt,
  left=3mm,right=3mm,top=2.3mm,bottom=2.3mm,fonttitle=\bfseries,coltitle=white}}
% --- boîtes TUTORIEL (noms EXACTS, ne pas renommer) ---
\newtcolorbox{cle}[1][]{boxbase,colback=defcol!6,colframe=defcol,
  title={\textsc{À retenir}\ifx\relax#1\relax\relax\else\textmd{ : \itshape #1}\fi}}
\newtcolorbox{methode}[1][]{boxbase,colback=methcol!7,colframe=methcol,sharp corners,
  title={\textsc{Méthode}\ifx\relax#1\relax\relax\else\textmd{ --- \itshape #1}\fi}}
\newtcolorbox{exemplebox}[1][]{boxbase,colback=excol!5,colframe=excol,
  title={\textsc{Exemple}\ifx\relax#1\relax\relax\else\textmd{ : \itshape #1}\fi}}
\newtcolorbox{piege}[1][]{boxbase,colback=attncol!6,colframe=attncol,
  title={\textsc{Piège}\ifx\relax#1\relax\relax\else\textmd{ : \itshape #1}\fi}}
\newtcolorbox{remarque}[1][]{boxbase,colback=black!4,colframe=black!45,
  title={\textsc{Remarque}\ifx\relax#1\relax\relax\else\textmd{ : \itshape #1}\fi}}
% --- boîtes EXERCICE / CORRIGÉ (noms EXACTS, auto-comptées) ---
\newtcolorbox[auto counter]{exo}[1][]{boxbase,colback=primary!4,colframe=primary,
  title={\textsc{Exercice}~\thetcbcounter\ifx\relax#1\relax\relax\else\textmd{ --- \itshape #1}\fi}}
\newtcolorbox[auto counter]{corrige}[1][]{boxbase,colback=excol!4,colframe=excol,
  title={\textsc{Corrigé}~\thetcbcounter\ifx\relax#1\relax\relax\else\textmd{ --- \itshape #1}\fi}}
\newcommand{\vv}[1]{\vec{#1}}\newcommand{\ds}{\displaystyle}
\newcommand{\R}{\mathbb{R}}\newcommand{\N}{\mathbb{N}}\newcommand{\Z}{\mathbb{Z}}
% =========================================================================
\begin{document}
\sisetup{per-mode=symbol}
% ================= CORRIGÉ DIFFICILE =================

\begin{corrige}[le satellite géostationnaire --- vrai ou faux]
\begin{enumerate}[label=\alph*)]
  \item \textbf{Vrai.} C'est la définition même du satellite géostationnaire.
  \item \textbf{Vrai.} Même période $T$ que la Terre $\Rightarrow$ même
        $\omega=2\pi/T$.
  \item \textbf{Faux.} $v=\omega r$ : le satellite est bien plus loin du centre
        ($r\gg R_T$), donc sa vitesse linéaire est \emph{bien plus grande} que celle
        de la maison.
  \item \textbf{Faux.} Ces conditions n'autorisent qu'une \emph{seule} altitude
        ($h\approx\qty{35800}{\kilo\metre}$).
\end{enumerate}
Correctes : a et b.
\end{corrige}

\begin{corrige}[classer trois satellites : $\omega$ et $v$]
\begin{enumerate}[label=\alph*)]
  \item $\omega=\sqrt{\mathcal{G}M_T/r^3}$ décroît quand $r$ augmente. Comme
        $r_3<r_2<r_1$ : $\boxed{\omega_3>\omega_2>\omega_1}$.
  \item $v=\sqrt{\mathcal{G}M_T/r}$ décroît aussi quand $r$ augmente :
        $\boxed{v_3>v_2>v_1}$.
\end{enumerate}
Le satellite le plus bas est le plus rapide, en rotation \emph{comme} en vitesse.
\end{corrige}

\begin{corrige}[à quelle altitude est-il géostationnaire ?]
\begin{enumerate}[label=\alph*)]
  \item On isole $r=\sqrt[3]{\dfrac{\mathcal{G}M_T\,T^2}{4\pi^2}}$ :
        \[
        r=\sqrt[3]{\frac{\num{4.0e14}\times\num{7.42e9}}{39{,}5}}
         =\sqrt[3]{\frac{\num{2.97e24}}{39{,}5}}
         =\sqrt[3]{\num{7.51e22}}\approx\qty{4.22e7}{\metre}.
        \]
  \item $h=r-R_T=\num{4.22e7}-\num{6.4e6}=\num{3.58e7}\,\si{\metre}
        =\qty{35800}{\kilo\metre}$. C'est l'unique altitude géostationnaire.
\end{enumerate}
\end{corrige}

\begin{corrige}[même $\omega$, mais quelle vitesse ?]
\begin{enumerate}[label=\alph*)]
  \item $v_{\text{sat}}=\omega\,r=\num{7.27e-5}\times\num{4.22e7}\approx\qty{3.1e3}{\metre\per\second}$.
  \item $v_{\text{maison}}=\omega\,R_T=\num{7.27e-5}\times\num{6.4e6}\approx\qty{465}{\metre\per\second}$.
  \item $\ds \frac{v_{\text{sat}}}{v_{\text{maison}}}=\frac{r}{R_T}
        =\frac{\num{4.22e7}}{\num{6.4e6}}\approx 6{,}6$ : le satellite va $\approx 6{,}6$
        fois plus vite. Même vitesse \emph{angulaire}, mais vitesses \emph{linéaires}
        très différentes, car $v=\omega r$ dépend du rayon.
\end{enumerate}
\end{corrige}

\end{document}
